\section{Simulation Results}
\label{sec:simulation_results}

Chapter~7 reports the comparative DTSE results under the frozen scenario grid defined in Chapter~6. The focus is on how the Onocoy anchor case and alternative mechanism profiles respond to matched stress inputs, and on the order in which those responses become visible across price, incentive, participation, and service fields.

\subsection{Reporting Conventions and Interpretation Boundaries}
\label{subsec:results_reporting_conventions}
This chapter reports the DePIN Tokenomic Stress Evaluator (DTSE) outputs for the stress scenarios defined in Section~\ref{sec:stress_scenarios}. Results are descriptive and comparative: each scenario evaluates how the Onocoy (ONO) anchor case and comparator DePIN mechanism profiles respond to matched stress inputs. As established in the methodology, these outputs represent model-conditional behavior under frozen assumptions; they are not empirical measurements, market forecasts, or absolute claims of real-world protocol superiority.

Chapter~5 also established that several ONO-relevant market and revenue fields remain not reliably observable in historical documentation alone. This results chapter therefore uses DTSE outputs to examine those fields under matched conditions rather than treating retrospective empirical observation as sufficient on its own.

All results in this chapter are governed by a strict sequence over magnitude reporting convention. Because selected metric families sit on different numeric scales across mechanism profiles, cross-profile outputs cannot be read as magnitude-equivalent rankings. The analytical focus is therefore on timing: which metric family registers stress first, how quickly that deviation propagates across subsystems, and whether service continuity degrades before or after participation erodes.

To operationalize this, scenario outcomes are evaluated entirely as deviations from their corresponding neutral baseline trajectories. Crucially, baseline trajectories can naturally drift over the 52-week reporting horizon due to deterministic mechanism rules. Stress is therefore defined and measured as a material acceleration, reversal, or threshold crossing relative to that drifting reference path, rather than against a static equilibrium state.

Signal timing follows the two-stage detection logic. Stage~1 identifies the earliest-trigger metric family, and Stage~2 classifies the subsequent pattern of deterioration into predefined operational failure-mode signatures. Governance adaptation, discretionary intervention, and other unmodeled human responses remain outside the evidentiary scope of this chapter and are taken up later in the discussion.

\subsection{Baseline Trajectories (The Neutral Reference)}
\label{subsec:baseline_trajectories}
Baseline runs establish the neutral reference trajectories for core DTSE state variables (e.g., modeled price signal, token supply and emissions/burns, provider participation, and utilization) in the absence of scenario-driven stress inputs. They serve as comparison anchors for this chapter and are not interpreted as static equilibrium states.

All baseline experiments use the standard DTSE configuration (as documented in Section~\ref{subsec:dtse_model_overview_assumptions}). In this neutral setup, we disable all stress scenario logic and hold exogenous inputs steady. This ensures that any baseline dynamics arise purely from the organic interaction between each protocol's tokenomic rule set and the modeled provider economics.

Baseline drift matters for interpretation. In DTSE, the baseline is the neutral reference path generated by the frozen rule set under neutral inputs; it is not a claim of long-term stability. Stress results are therefore strictly compared against that drifting reference path rather than against a flat line.

In the ONO baseline, for example, modeled price and active provider counts naturally decline over the 52-week horizon due to expected token emission decay. Under stress, the diagnostic question is not whether these metrics fall, but whether they fall faster, reverse earlier, or cross critical thresholds sooner than they would have under neutral conditions. This is how the chapter reads every subsequent scenario: by measuring when and where a stress input pulls the system away from its own baseline trajectory, and which metric family registers the departure first.

\begin{figure}[H]
    \centering
    \includegraphics[width=0.95\linewidth]{../exports/dtse/figures/fig_baseline_ono_price_providers.png}
    
    \vspace{0.4\baselineskip}
    
    \includegraphics[width=0.95\linewidth]{../exports/dtse/figures/fig_baseline_ono_burn_mint.png}
    \caption[DTSE ONO baseline reference trajectories]{DTSE ONO baseline reference trajectories under neutral inputs: modeled price and active provider count (top) and emissions versus sinks (bottom). These curves provide the drifting reference path against which later stress deviations are read.}
    \label{fig:dtse_baseline_reference}
\end{figure}

\subsection{Cross-Scenario Detection Pattern Summary}
\label{subsec:mechanism_agnostic_response}
Across the DTSE experiment set, each stress channel produces a distinct sequence of metric changes when read as baseline deviations. The purpose of this section is to summarize those recurring response patterns before the chapter turns to the ONO-specific scenario figures. Table~\ref{tab:dtse_mechanism_agnostic_framework} maps each modeled stress channel to its earliest signal, its participation-side follow-on field, and the Stage~2 signature observed under that response pattern.

\begin{table}[ht]
    \centering
    \footnotesize
    \setlength{\tabcolsep}{3pt}
    \renewcommand{\arraystretch}{1.3}
    \begin{tabularx}{\textwidth}{@{}
      >{\raggedright\arraybackslash}p{0.24\textwidth}
      >{\raggedright\arraybackslash}p{0.20\textwidth}
      >{\raggedright\arraybackslash}p{0.22\textwidth}
      >{\raggedright\arraybackslash}X
    @{}}
    \toprule
    \textbf{Stress Channel (Modeled Input)} & \textbf{Earliest Signal} & \textbf{Participation-Side Follow-On} & \textbf{Observed Stage~2 Signature} \\
    \midrule
    \textbf{Demand Contraction:}\\Exogenous service-usage reduction under a specified demand regime. & Utilization; usage-linked sinks. & Active provider count deteriorates materially after usage-linked weakening. & \textbf{Reward--Demand Decoupling} \\
    \addlinespace
    \textbf{Liquidity Shock:}\\Discrete price dislocation via modeled finite-depth liquidity channel. & Modeled price signal. & Provider profitability and provider churn deteriorate around the event window. & \textbf{Liquidity-Driven Compression} \\
    \addlinespace
    \textbf{Competitive Yield Pressure:}\\Exogenous increase in outside-opportunity attractiveness. & Provider churn rate. & Active provider count weakens before utilization becomes the first-moving field. & \textbf{Elastic Provider Exit} \\
    \addlinespace
    \textbf{Provider Cost Inflation:}\\Exogenous upward shift in operator hardware/OpEx costs. & Provider profitability. & Active provider count and provider churn weaken after sustained margin compression. & \textbf{Profitability-Induced Churn} \& \textbf{Latent Capacity Degradation} \\
    \bottomrule
    \end{tabularx}
    \caption[Cross-scenario detection pattern summary]{Cross-scenario detection pattern summary under the frozen DTSE setup. The table reports how each modeled stress channel first appears and which Stage~2 signature is observed once deterioration propagates across subsystems.}
    \label{tab:dtse_mechanism_agnostic_framework}
\end{table}

\subsection{Onocoy (ONO) Baseline-Relative Signal Sequencing}
\label{subsec:ono_signal_sequencing}
Table~\ref{tab:dtse_ono_signal_sequencing} makes the same sequencing logic concrete for the Onocoy (ONO) profile under the frozen stress scenarios. The table is read strictly against the ONO baseline: it shows which subsystems move first, which fields follow, and which Stage~2 signature is observed under each channel without claiming exact live-network thresholds or token-price outcomes.

\begin{table}[ht]
    \centering
    \footnotesize
    \setlength{\tabcolsep}{3pt}
    \renewcommand{\arraystretch}{1.3}
    \begin{tabularx}{\textwidth}{@{}
      >{\raggedright\arraybackslash}p{0.18\textwidth}
      >{\raggedright\arraybackslash}p{0.25\textwidth}
      >{\raggedright\arraybackslash}p{0.25\textwidth}
      >{\raggedright\arraybackslash}X
    @{}}
    \toprule
    \textbf{Stress Channel} & \textbf{First-Moving Field (Stage 1 Detection)} & \textbf{Secondary Detection Field (Stage 2)} & \textbf{Diagnosed Signature for ONO} \\
    \midrule
    \textbf{Demand Contraction} & Utilization and usage-linked sinks deviate first. & Active provider count deviates materially later. & Reward--Demand Decoupling. \\
    \addlinespace
    \textbf{Liquidity Shock} & Modeled price drawdown triggers immediately. & Provider profitability and provider churn break following the price event. & Liquidity-Driven Compression. \\
    \addlinespace
    \textbf{Competitive Yield Pressure} & Provider churn rate and active provider count deviate immediately. & Utilization remains stable initially, lagging behind participation drop. & Elastic Provider Exit. \\
    \addlinespace
    \textbf{Provider Cost Inflation} & Provider profitability turns structurally negative. & Active provider count and provider churn degrade following margin collapse. & Profitability-Induced Churn \& Latent Capacity Degradation. \\
    \bottomrule
    \end{tabularx}
    \caption[Onocoy (ONO) baseline-relative signal sequencing]{Onocoy (ONO) baseline-relative signal sequencing under the frozen stress scenarios. The table summarizes which fields move first, which fields follow, and which Stage~2 signature is observed in each case.}
    \label{tab:dtse_ono_signal_sequencing}
\end{table}

\subsection{Scenario-by-Scenario Deviations}
\label{subsec:scenario_by_scenario_deviations}
The preceding tables summarized each stress channel's detection sequence in compact form. This section provides the detailed simulation evidence behind those summaries, walking through the observed scenario deviations for the ONO baseline one channel at a time. The focus remains on which metric families register first and how the initial deviation propagates through later fields.

\subsubsection{Demand Contraction (Reward--Demand Decoupling)}
Demand contraction isolates a scenario where real-world usage weakens while reward logic remains active. The diagnostic question is which subsystem registers first under baseline-relative detection, and whether participation effects appear only after profitability and provider churn thresholds begin to bind.

Under this channel, utilization and usage-linked sink fields deviate from the baseline first. However, the active provider count does not meaningfully break from the baseline trajectory until detectably later in the simulation horizon. The network continues its scheduled reward distribution while users are no longer buying sufficient Data Credits to offset the outflow.

Taken together, the ONO demand-contraction run shows a clear \emph{Reward--Demand Decoupling} pattern. Usage-linked fields weaken first, while provider participation remains comparatively sticky and only deteriorates materially much later. The result is not immediate network discontinuity, but a widening gap between ongoing reward distribution and weakening service-side demand.

\begin{figure}[H]
    \centering
    \includegraphics[width=0.98\linewidth]{../exports/dtse/figures/fig_demand_contraction_ono_core.png}
    \caption[DTSE ONO demand-contraction deviations from baseline]{DTSE ONO demand-contraction deviations from baseline: utilization, sinks, and provider profit. The channel breaks first in usage-linked fields, with participation effects arriving materially later.}
    \label{fig:dtse_demand_contraction_ono}
\end{figure}

\subsubsection{Liquidity Shock (Liquidity-Driven Compression)}
The liquidity shock scenario isolates the transmission of a discrete token-liquidity dislocation into provider-side stress fields. The simulation applies a discrete token sell-off directly against a finite-depth liquidity channel, creating an immediate price drop without an accompanying initial collapse in real-world service demand.

When the configured token unlock event hits the market, the abrupt price dislocation immediately compresses the fiat-equivalent reward value for providers. Within a short window, provider profitability falls below zero and provider churn rises sharply. These signals appear together even though the scenario does not begin with a simultaneous collapse in service demand.

In the ONO liquidity-shock run, the decisive feature is the short transmission distance between the price event and provider-side deterioration. The modeled price dislocation compresses fiat-equivalent reward value around the event window, and profitability and churn respond almost immediately thereafter. The observed Stage~2 pattern is therefore \emph{Liquidity-Driven Compression} rather than a demand-led contraction.

\begin{figure}[H]
    \centering
    \includegraphics[width=0.98\linewidth]{../exports/dtse/figures/fig_liquidity_shock_ono_price_churn.png}
    \caption[DTSE ONO liquidity-shock transmission]{DTSE ONO liquidity-shock transmission: modeled price dislocation and provider churn response. The figure highlights how a financial shock is transmitted into participation-side stress around the event window.}
    \label{fig:dtse_liquidity_shock_ono}
\end{figure}

\subsubsection{Competitive Yield Pressure (Elastic Provider Exit)}
This channel isolates provider sensitivity to outside opportunities without requiring internal demand collapse or a discrete internal liquidity event. In this scenario, operators suddenly gain access to a much more attractive alternative for their hardware or capital, parameterized to reflect competing yield networks like adjacent hardware protocols or re-purposed GNSS arrays.

As soon as the outside yield pressure increases, participation fields register immediate deviations. Provider churn rises, and the active provider count drops sharply in the earliest stages of the simulation. During this initial period, the internal network utilization metrics remain comparatively stable, so the first visible change is on the supply side rather than the service side.

For ONO, competitive-yield pressure first appears as a participation-side response to a more attractive outside option. Churn and active-provider counts deviate before utilization becomes the first-moving field, indicating that operator withdrawal can begin while service-side conditions still appear comparatively stable. The observed Stage~2 pattern is therefore \emph{Elastic Provider Exit}.

\begin{figure}[H]
    \centering
    \includegraphics[width=0.98\linewidth]{../exports/dtse/figures/fig_competitive_yield_ono_retention_churn.png}
    \caption[DTSE ONO competitive-yield pressure response]{DTSE ONO competitive-yield pressure response: active provider count and provider churn deviations versus baseline. Participation moves before visible service-side deterioration when outside yield becomes more attractive.}
    \label{fig:dtse_competitive_yield_ono}
\end{figure}

\subsubsection{Provider Cost Inflation (Latent Capacity Degradation)}
Provider cost inflation tests exposure to real-world operating environments (e.g., electricity prices, connectivity costs, facility rents). The scenario applies an exogenous upward shift to the physical operating-cost conditions of the providers.

This channel immediately triggers the \emph{Profitability-Induced Churn} signature: rising fiat-equivalent costs compress provider margins into negative territory immediately. However, the simulation also shows that the active provider count does not synchronously collapse. Active nodes begin to deteriorate only after sustained windows of unprofitability.

For ONO, provider-cost inflation first appears as margin compression and only later as visible participation loss. The run therefore shows a two-step deterioration pattern: profitability weakens immediately, while active-provider deterioration follows after sustained unprofitability. In Stage~2 terms, the scenario expresses \emph{Profitability-Induced Churn} together with \emph{Latent Capacity Degradation}.

\begin{figure}[H]
    \centering
    \includegraphics[width=0.9\linewidth]{../exports/dtse/figures/fig_cost_inflation_ono_profit_providers.png}
    \caption[DTSE ONO provider-cost inflation response]{DTSE ONO provider-cost inflation response: margin compression and subsequent active provider count dynamics. The figure shows profitability deteriorating before participation metrics visibly break from baseline.}
    \label{fig:dtse_cost_inflation_ono}
\end{figure}

\begin{table}[ht]
    \centering
    \footnotesize
    \setlength{\tabcolsep}{3pt}
    \renewcommand{\arraystretch}{1.25}
    \begin{tabularx}{\textwidth}{@{}
      >{\raggedright\arraybackslash}p{0.18\textwidth}
      >{\raggedright\arraybackslash}p{0.25\textwidth}
      >{\raggedright\arraybackslash}p{0.28\textwidth}
      >{\raggedright\arraybackslash}X
    @{}}
    \toprule
    \textbf{Stress Channel} & \textbf{Earliest Material Deviation} & \textbf{Observed Propagation Pattern} & \textbf{Observed Stage~2 Signature} \\
    \midrule
    Demand Contraction & Utilization and usage-linked sinks weaken before provider participation does. & Reward distribution remains active while service-side demand weakens; participation only deteriorates later. & Reward--Demand Decoupling \\
    \addlinespace
    Liquidity Shock & Modeled price dislocation appears around the event window. & Price compression is followed quickly by weaker provider economics and higher churn. & Liquidity-Driven Compression \\
    \addlinespace
    Competitive Yield Pressure & Participation fields break before utilization becomes the first-moving field. & Outside-option pressure appears first as churn and active-provider deterioration while service-side measures remain comparatively stable. & Elastic Provider Exit \\
    \addlinespace
    Provider Cost Inflation & Provider profitability turns negative before participation breaks materially. & Margin compression persists first, then propagates into provider churn and active-provider deterioration. & Profitability-Induced Churn + Latent Capacity Degradation \\
    \bottomrule
    \end{tabularx}
    \caption[Observed Stage~2 summary for ONO]{Observed Stage~2 summary for the ONO profile under the frozen DTSE scenarios. The table closes the scenario readouts by showing how earliest deviations propagate into diagnosed failure signatures.}
    \label{tab:dtse_ono_stage2_summary}
\end{table}

\subsection{Cross-Profile and Cross-Scenario Synthesis}
\label{subsec:dtse_cross_profile_synthesis}
Having diagnosed the operational failure signatures for each stress channel individually, this section synthesizes the results across all four scenarios and compares them against alternative mechanism profiles.

\begin{figure}[H]
    \centering
    \includegraphics[width=0.85\linewidth]{../exports/dtse/figures/fig_cross_scenario_ono_metric_deltas_heatmap.png}

    \vspace{0.2\baselineskip}

    \includegraphics[width=0.85\linewidth]{../exports/dtse/figures/fig_cross_scenario_signal_timing_ono.png}
    \caption[DTSE ONO cross-scenario synthesis]{DTSE ONO cross-scenario synthesis: final-metric delta heatmap (top) and first-signal timing by stress channel (bottom). Together, the panels demonstrate both where stress concentrates operationally and how quickly each channel becomes visible in the pre-specified metric suite.}
    \label{fig:dtse_signal_timing_chart}
\end{figure}

\noindent\textbf{Cross-Scenario Observation:} Figure~\ref{fig:dtse_signal_timing_chart} makes the sequence-over-magnitude pattern visible across the ONO stress set. The heatmap and timing charts show that stress rarely strikes the physical participation layer first. Across nearly all channels, financial indicators (profitability) or usage indicators (utilization) act as the early-warning layer, while active provider counts tend to move later.

Across the experiment set, stress channels manifest in distinct metric families and timing windows (Figures~\ref{fig:dtse_cross_profile_retention} and~\ref{fig:dtse_signal_timing_chart}). Demand-side stress is detected earlier in utilization-linked channels. Liquidity-shock signals appear in price, provider-economics, and provider-churn fields around the event window. Cost and yield channels are first detected in provider-economics and provider-churn fields.

\begin{figure}[H]
    \centering
    \includegraphics[width=0.85\linewidth]{../exports/dtse/figures/fig_cross_scenario_retention_profiles.png}
    
    \vspace{0.2\baselineskip}
    
    \includegraphics[width=0.85\linewidth]{../exports/dtse/figures/fig_cross_scenario_solvency_profiles.png}
    \caption[DTSE cross-profile final-week comparison by scenario]{DTSE cross-profile final-week comparison by scenario: provider retention (top) and incentive-solvency proxy (bottom). These outputs must be read directionally under matched stress inputs rather than as magnitude-equivalent live-protocol rankings.}
    \label{fig:dtse_cross_profile_retention}
\end{figure}

\noindent\textbf{Cross-Profile Observation:} Figure~\ref{fig:dtse_cross_profile_retention} shows that mechanism profiles channel matched stress differently. Under identical liquidity and yield shocks, BME-oriented profiles exhibit larger swings in their incentive-solvency proxies. In contrast, the capped-supply ONO profile absorbs the same pressures through a more bounded proxy path, with stress appearing more clearly in provider margin compression.

\noindent\textbf{Incentive-Solvency Proxy Caveat:} The proxy used in Figure~\ref{fig:dtse_cross_profile_retention} serves as a comparative stress signal rather than a standalone economic health score. For capped-supply profiles like ONO, the ``mint'' denominator represents scheduled reward distributions from fixed inventory rather than uncapped token creation. A larger proxy swing (such as those seen in BME comparators) therefore indicates a different structural transmission path under stress, not a categorically ``worse'' absolute outcome.

Because the proxy behaves like a ratio, it can move sharply when one side deteriorates faster than the other. Accordingly, the proxy bars in Figure~\ref{fig:dtse_cross_profile_retention} are interpreted jointly with profitability, provider churn, provider retention, and utilization rather than in isolation.

In summary, the cross-scenario synthesis supports a sequential reading of DePIN stress under the frozen DTSE setup. Financial, usage-linked, and provider-economics fields absorb the initial shock first, while physical participation metrics often move later. These observed sequences provide the basis for the discussion that follows.
